\documentclass[11pt, letterpaper]{article}
\input{/home/hyperion/.config/LatexHub/preambles/base.tex}
\input{/home/hyperion/.config/LatexHub/preambles/diagrams.tex}
\input{/home/hyperion/.config/LatexHub/preambles/tikz.tex}
\input{/home/hyperion/.config/LatexHub/preambles/macros-cat-theory.tex}
\input{/home/hyperion/.config/LatexHub/preambles/macros-algebra.tex}
\input{/home/hyperion/.config/LatexHub/preambles/basic-theorems-section.tex}
\input{/home/hyperion/.config/LatexHub/preambles/refs-basic.tex}
\theoremstyle{plain}
\newtheorem{problem}{Problem}
\usepackage{titlepageBU}
\title{Homework 8}
\courseID{MA731}
\begin{document}
\makeproblem


\begin{problem}
	Let $\mathfrak{g} $ be a 5-dimensional Lie algebra over a field of characteristic 0. Suppose that with respect to a basis $\left\{ e_1, \ldots , e_5 \right\} $, every $\mathrm{ad} _x$ is strictly upper triangular, and the pattern of nonzero entries is given by
	\[
		\begin{pmatrix}
		0 & 0 & * & * & * \\
		0 & 0 & 0 & * & * \\
		0 & 0 & 0 & 0 & * \\
		0 & 0 & 0 & 0 & * \\
		0 & 0 & 0 & 0 & 0
		\end{pmatrix}
	\]
	where each ``$*$'' indicates an entry that is nonzero (for suitable $x\in \mathfrak{g} $), and all other entries are zero.

The \textbf{upper central series} is defined inductively by
\[
	Z_0 = 0,\qquad Z_{i+1} =\left\{ x\in \mathfrak{g} \mid [x,\mathfrak{g} ]\subseteq Z_i \right\} 
.\]
The \textbf{lower central series} is defined by
\[
	\mathfrak{g} _1 = \mathfrak{g} ,\qquad \mathfrak{g} _{i+1} =[\mathfrak{g} ,\mathfrak{g} _i]
.\]
\begin{enumerate}
	\item Express each $Z_i$ as the span of a subset of the basis vectors $\left\{ e_1, \ldots ,e_5 \right\} $.
	\item Express each $\mathfrak{g} _i$ as the span of a subset of the basis vectors $\left\{ e_1, \ldots ,e_5 \right\} $.
	\item Give an example of a nilpotent Lie algebra (not necessarily of dimension 5) for which the sequence $\dim Z_i$ is \textbf{not} the reverse of the sequence $\dim \mathfrak{g} _i$.
\end{enumerate}
\end{problem}
\begin{proof}
	(1,2) We adopt the following notation: for each $x\in \mathfrak{g}\setminus \left\{ e_5 \right\}  $, write
	\[
		\mathrm{ad} _x = \begin{pmatrix}
		0 & 0 & x_{13}  & x_{14}  & x_{15}  \\
		0 & 0 & 0 & x_{24}  & x_{25}  \\
		0 & 0 & 0 & 0 & x_{35}  \\
		0 & 0 & 0 & 0 & x_{45}  \\
		0 & 0 & 0 & 0 & 0
		\end{pmatrix}\quad \text{ and } \quad 
		\mathrm{ad} _{e_5} =\begin{pmatrix}
		0 & 0 & a_{13}  & a_{14}  & a_{15}  \\
		0 & 0 & 0 & a_{24}  & a_{25}  \\
		0 & 0 & 0 & 0 & a_{35}  \\
		0 & 0 & 0 & 0 & a_{45}  \\
		0 & 0 & 0 & 0 & 0
		\end{pmatrix}
	.\]
	The assumption that each ``$*$'' is nonzero for suitable $x\in \mathfrak{g} $ in particular implies that for each $*$, there is at least one $x\in \mathfrak{g} $ for which that component of $\mathrm{ad} _x$ is nonzer. In particular, there exists $x\in \mathfrak{g} $ such that $x_{35} \neq 0$ or $x_{45} \neq 0$.

	Let $x\in \mathfrak{g} $. Since $\left\{ e_1, \ldots ,e_5 \right\} $ is a basis, write $x = \sum_{i} c_ie_i$ for coefficients $c_i\in\mathbb{K} $, with $\mathbb{K} $ the base field. Then we have
	\[
		[e_5,x] = \sum_{1\leq i\leq 5} c_i [e_5,e_i]
	.\]
	Note that for any $y\in \mathfrak{g} $, the structure of the matrix $\mathrm{ad} _y$ implies $\mathrm{ad} _y(e_1) = \mathrm{ad} _y(e_2) = 0$. Also note that $[e_5,e_5]=0$. Thus taking $y=e_5$ simplifies the above sum to
	\[
		[e_5,x] = c_3[e_5,e_3] + c_4[e_5,e_4] = c_3\mathrm{ad}_{e_5} (e_3) + c_4 \mathrm{ad} _{e_5} (e_4)
	.\]
	Note that $\mathrm{ad} _{e_5} (e_3) = a_{13} e_1$ and $\mathrm{ad} _{e_5} (e_4) = a_{14} e_1 + a_{24} e_2$. In particular, $\mathrm{ad} _{e_5} (e_3) + \mathrm{ad} _{e_5} (e_4)$ is linearly independent from $\left\{ e_3,e_4 \right\} $, and thus contains no $e_3$- or $e_4$-dependence. Thus $[e_5,x] = -[x,e_5]$ contains no $e_3$- or $e_4$-dependence. Therefore
	\[
		-[x,e_5] = -\mathrm{ad} _x(e_5) = -\left( x_{15} e_1 + x_{25} e_2 + x_{35} e_3 + x_{45} e_4 \right) 
	\]
	must not depend on $e_3$ or $e_4$, implying $-x_{35} =-x_{45} =0$. Thus $x_{35} =x_{45} =0$. Since $x\in \mathfrak{g} $ was arbitrary, for all $x\in \mathfrak{g} $, we have $x_{35} =x_{45} =0$, contradicting the assumptions of the problem. Either the problem assumes a false statement is true, implying that all statements are true (including the claims of the problem), or $\mathfrak{g} $ does not exist, implying (1) and (2) vacuously. We proceed with the second assumption, and prove (3).

	(3) Let $\mathfrak{g} $ a 1-dimensional Lie algebra and $\mathfrak{h} $ the Heisenberg algebra generated by the elements
	\[
		X=\begin{pmatrix}
		0 & 1 & 0 \\
		0 & 0 & 0 \\
		0 & 0 & 0
		\end{pmatrix},\qquad Y = \begin{pmatrix}
		0 & 0 & 0 \\
		0 & 0 & 1 \\
		0 & 0 & 0
		\end{pmatrix},\qquad Z = \begin{pmatrix}
		0 & 0 & 1 \\
		0 & 0 & 0 \\
		0 & 0 & 0
		\end{pmatrix}
	.\]
	We show that $\mathfrak{k} \coloneqq \mathfrak{g} \oplus \mathfrak{h} $ is an example of (3). Since $\mathfrak{h} $ has dimension 3, $\mathfrak{k} $ has dimension 4. Let $e\in \mathfrak{g} $ be the basis, such that $\mathfrak{k} $ has the basis
	\[
		\left\{ e,X,Y,Z \right\} 
	.\]
	Here we note that $[e,\mathfrak{h} ]=0$ (since we are identifying $e$ with the ordered pair $(e,0)$, and similarly for the elements of $\mathfrak{h} $). We also recall the commutation relations of $\mathfrak{h} $:
	\[
		[X,Y]=Z,\qquad [X,Z] =0,\qquad  [Y,Z] = 0
	.\]
	Since the only nonzero commutator between basis elements is $[X,Y]=Z$, we immediately see that $\mathfrak{g} _2 = [\mathfrak{g} ,\mathfrak{g} ] = \operatorname{span} \left\{ Z \right\} $ has dimension 1. Since the lower central series is indexed starting at 1, and the upper central series is indexed starting at 2, it suffices to show $\dim Z_1 \neq 1$.

	Recall $Z_1 = \left\{ x\in \mathfrak{k} \mid \mathrm{ad} _x=0 \right\} $. We immediately see that $\mathfrak{g} $ must be abelian since $[e,e]=0$, so we immediately have $\mathrm{ad} _e=0$. Additionally, we have $\mathrm{ad} _Z(X) =\mathrm{ad} _Z(Y) = 0$ and $\mathrm{ad} _Z(e)=0$, implying $\mathrm{ad} _Z =0$. Thus we have at least two linearly independent elements $e,Z\in Z_1$, implying $\dim Z_1 \geq 2 \neq 1$. Therefore $\mathfrak{k} $ is an example of such an algebra.
\end{proof}

\begin{proof}
	I also solve (1,2) in the intended way. For all $i$, note that $\mathfrak{g} _i$ and $Z_i$ are all clearly subspaces of $\mathfrak{g} $, and note that we trivially have $Z_i \subseteq Z_{i+1} $ and $\mathfrak{g} _i \supseteq \mathfrak{g} _{i+1} $.

	(1) Let $x\in \mathfrak{g} $; with notation as above,
	\[
		\mathrm{ad}_x(e_1) = 0,\qquad \mathrm{ad}_x(e_2)=0,
	\]
	\[
		\mathrm{ad} _x(e_3) = x_{13} e_1,\qquad \mathrm{ad} _x(e_4)=x_{14} e_1 + x_{24} e_2,
	\]
	\[
		\mathrm{ad} _x(e_5) = x_{15} e_1 + x_{25} e_2 + x_{35} e_3 + x_{45} e_4
	.\]
	We assume that for all ``$*$'' in the given form, there exists $x$ such that that component is nonzero. Note that $Z_1 = \left\{ x\in \mathfrak{g} \mid \mathrm{ad}_x=0 \right\} $. Linear independence of $e_1,e_2,e_3,e_4$ and the above equations imply $Z_1 = \operatorname{span} \left\{ e_1,e_2 \right\} $. Then $Z_2$ is the set of all $x$ for which $\mathrm{ad} _x(y)\in \operatorname{span} \left\{ e_1,e_2 \right\} $ for all $y\in \mathfrak{g} $. The above equations imply this is $Z_2 = \left\{ e_1,e_2,e_3,e_4 \right\} $. By the same logic, we see $Z_3 = \operatorname{span} \left\{ e_1, \ldots ,e_5 \right\} $.

	(2) Note that $\mathfrak{g} _2 = [\mathfrak{g} ,\mathfrak{g} ]$. Letting $x\in \mathfrak{g} $ such that $x_{13} \neq 0$, and using that $\mathfrak{g} _2$ is a subspace, we first see that $x_{13} ^{-1} \mathrm{ad} _x(e_3) = e_1\in \mathfrak{g} _2$. Now let $x$ have $x_{24} \neq 0$. We have
	\[
		x_{24} ^{-1} \left(\mathrm{ad} _x(e_4) - x_{14} e_1 \right)= e_2\in \mathfrak{g} _2,
	\]
	since $e_1\in \mathfrak{g}_2$ and thus $\mathfrak{g} _2$ is closed under subtaction by it. From this we similarly obtain $x_{35} e_3 + x_{45} e_4\in \mathfrak{g}_2$. We cannot obtain $e_5\in \mathfrak{g} _2$, so $\mathfrak{g} _2$ has dimension at most 4. I cannot find how to conclude that $e_3,e_4\in \mathfrak{g} _2$, but the statement of (3) implies that this must be true. If one assumes that $x_{35} $ and $x_{45} $ are independent parameters this follows, but the problem does not state that. We proceed with the assumption that $\mathfrak{g} _2 = \operatorname{span} \left\{ e_1, \ldots ,e_4 \right\} $, but note that if one takes $\mathfrak{g} _2 = \operatorname{span} \left\{ e_1, e_2, e_3 + ce_4 \right\} $ for some $c\in\mathbb{K} $, then the logic below does not change by bilinearity of the bracket.

	Now we immediately see that $\mathfrak{g} _3 = [x,y]$ with $x\in \mathfrak{g} $ and $y\in \operatorname{span} \left\{ e_1, \ldots ,e_4 \right\} $. Recall that $\mathrm{ad} _{e_1} =\mathrm{ad} _{e_2} =0$, so we get no information there. We then have that there is $x$ such that $[x,e_3]=x_{13} e_1$, so $e_1\in \mathfrak{g} _3$. We then obtain that $e_2\in \mathfrak{g} _3$ by $e_2 = x_{24} ^{-1} [x,e_4] - x_{24} ^{-1} x_{14} e_1$ for suitable $x$. Additionally, 
	\[
		[e_5,e_3] = \mathrm{ad}_{e_5} (e_3)\in \operatorname{span} \left\{ e_1,e_2 \right\} \ni \mathrm{ad}_{e_5} (e_4) = [e_5,e_4]
	.\]
	Thus $\mathfrak{g} _3 = \operatorname{span} \left\{ e_1,e_2 \right\} $, since we have exhausted all combinations of basis elements. Finally, since $\mathrm{ad} _{e_1} =\mathrm{ad} _{e_2} =0$, we obtain $\mathfrak{g} _4=0$.
\end{proof}





\end{document}
